\section{非前向与遍举量：准 GPD、准 DA 及其他}\label{sec:lamet-offforward}

同样的构造适用于非前向矩阵元与真空到强子的矩阵元。GPD 与 DA 的光锥定义见\cref{def:excl-gpd,def:excl-da}；本书的偏斜度 $\xi=\frac{p^+-p'^+}{p^++p'^+}$（Ji 的 $\xi_{\rm Ji}=2\xi$）。

\begin{definition}[准 GPD 与准 DA]\label{def:lamet-quasi-gpd}
\begin{enumerate}[label=(\roman*),leftmargin=2em,itemsep=1pt,topsep=2pt]
\item \emph{准 GPD}（对称参考系：入射 $P$、出射 $P'$，$\bar P=\frac{P+P'}2=(\bar P^0,0,0,\bar P^z)$，$\Delta=P'-P$，$t=\Delta^2$，$\tilde\xi=-\Delta^z/(2\bar P^z)$）：对 $\Gamma=\gamma^\alpha$，$\alpha\in\{0,z\}$，
\begin{equation}\label{eq:lamet-quasi-gpd}
\begin{split}
&\int\frac{\dd z}{4\pi}\,\ee^{\ii y\bar P^zz}\,\langle P'|\big[\bar\psi(\tfrac z2\hat e_z)\,\Gamma\, W[\tfrac z2\hat e_z,-\tfrac z2\hat e_z]\,\psi(-\tfrac z2\hat e_z)\big]_R|P\rangle\\
&\qquad=\frac1{2\bar P^z}\Big[\tilde H_\Gamma\,\bar u(P')\Gamma u(P)+\tilde E_\Gamma\,\bar u(P')\frac{\ii\sigma^{\alpha\nu}\Delta_\nu}{2M}u(P)\Big],
\end{split}
\end{equation}
$\tilde H_\Gamma=\tilde H_\Gamma(y,\tilde\xi,t,\bar P^z)$ 等。前向极限 $\tilde H_\Gamma(y,0,0,P^z)=\tilde q_\Gamma(y,P^z)$。对称系中 $\bar P\cdot\Delta=0$，故 $\xi=-\Delta^+/(2\bar P^+)=\tilde\xi\,\bar P^z/\bar P^0\to\tilde\xi$（$\bar P^z\to\infty$）。
\item \emph{准 DA}：$\langle0|\bar d(z\hat e_z)\gamma^z\gamma_5W[z\hat e_z,0]u(0)|\pi^+(P)\rangle_R=\ii f_\pi P^z\int\dd x\,\ee^{\ii(1-x)P^zz}\tilde\phi(x,P^z)$，$x$ 为 $u$ 夸克的动量分数；由定域极限 $\langle0|\bar d\gamma^\mu\gamma_5u|\pi^+(P)\rangle=\ii f_\pi P^\mu$，$\int\dd x\,\tilde\phi=1$。
\end{enumerate}
\uses{def:lattice-quasi, def:excl-gpd, def:excl-da}
\end{definition}

前向极限的检验：$P'=P$ 时把算符平移 $z/2$，左边化为 $\tilde q_\Gamma/N_\Gamma$；右边为 $\frac1{2P^z}\tilde H_\Gamma\cdot2P^\Gamma$；由 $N_\Gamma=P^z/P^\Gamma$ 得 $\tilde H_\Gamma=\tilde q_\Gamma$。准 DA 的相位：把\cref{def:excl-da}中以 $\pm z$ 为端点的算符平移 $z$，相位变为 $\ee^{-\ii(1-x)p\cdot(2z)}$；对类空间隔 $p\cdot z=-p^zz$，这就是 $\ee^{\ii(1-x)P^zz}$；对 $\ee^{\ii(1-x)P^zz}$ 积分 $\int\frac{\dd z}{2\pi}$ 给出 $\delta((1-x)P^z)$，由此得 $\int\tilde\phi=1$。不对称参考系与与参考系无关的 Lorentz 振幅定义见\cite{Bhattacharya:2022aob}（据其摘要，后者“可能给出更小的幂次修正”）；本节的参考系约定同\cite{Su:2026eta}（预印本）。

\begin{theorem}[准 GPD 与准 DA 的匹配 \tagA]\label{thm:lamet-quasi-gpd}
准 GPD 与准 DA 满足
\begin{equation}\label{eq:lamet-gpd-match}
\begin{aligned}
\tilde H(y,\xi,t,\bar P^z)&=\int_{-1}^1\frac{\dd x}{|x|}\,\tilde C\Big(\frac yx,\frac\xi x,\frac\mu{|x|\bar P^z}\Big)H(x,\xi,t,\mu)+(\text{幂次修正}),\\
\tilde\phi(x,P^z)&=\int_0^1\dd y\,C\Big(x,y,\frac\mu{P^z}\Big)\phi(y,\mu)+(\text{幂次修正}),
\end{aligned}
\end{equation}
$\tilde E$ 同样；$\tilde C$ 与 $t$ 无关（至少在单圈），在 $\xi=0$ 处等于 PDF 的核。LaMET 预言在 $[-1+x_0,-\xi-x_0]\cup[-\xi+x_0,\xi-x_0]\cup[\xi+x_0,1-x_0]$ 上可靠，$x_0$ 由硬部分子动量决定\cite{Holligan:2025baj}（阈值研究，据其摘要）。（本定理只涉及部分子态；强子层次的处理与前向情形（\cref{sec:lamet-power}）相同，未对 GPD 单独表述条件定理。）
\uses{def:lamet-quasi-gpd, asm:lamet-ope, prop:lattice-renormalizability, prop:lamet-support}
\end{theorem}

\begin{proofsketch}
OPE 路线：(1) 乘法重整化（\cref{prop:lattice-renormalizability}）；(2) 非前向的\cref{asm:lamet-ope}，算符为共形扭度 2 算符，其非前向矩阵元是 $H$、$E$ 的矩（多项式性）；(3) 支撑与解析性的论证（\cref{prop:lamet-support}）推广到非前向运动学，重求和给出坐标空间匹配，再作 Fourier 变换；部分子动量的论证与前向情形相同。单圈（均据其摘要）：非极化与螺旋度\cite{Ji:2015qla}（单圈领头幂次下 $E$（$\tilde E$）“不需要匹配”），横度\cite{Xiong:2015nua}，RI/MOM\cite{Liu:2019urm}，包括胶子与任意自旋的全部准 GPD\cite{Ma:2022ggj}，坐标空间\cite{Radyushkin:2019owq}，统一的 NLO\cite{Yao:2022vtp}；介子 DA 的单圈见 GPD 文献，RI/MOM\cite{Liu:2018tox}（准 PDF 的重整化因子可以使用，“只要在格点与微扰论中以同样方式实施”），格点\cite{Zhang:2017bzy,LatticeParton:2022zqc,Holligan:2023rex}；非单态 GPD 与介子 DA 的 NNLO 坐标核\cite{Ji:2025mvk}（预印本，单一小组）。完整论证见\cite{JiEtAl2021}及其所引文献。
\end{proofsketch}

\paragraph{地位与缺口}
\tagA；瓶颈是非前向形式的\cref{asm:lamet-ope}。保证平移不变性所需的运动学扭度 3 贡献按对数增强，并在 $x=\pm\xi$ 处不连续\cite{Braun:2023alc}（据其摘要）；$x=\pm\xi$ 这一点的状况见\cref{conj:yzliu-gpd-xi}（第~23~章）。没有 \tagR 条目：本书没有独立核对其中任何定阶结果（以上引述均据各文摘要）。

\begin{proposition}[其他准分布：重子与重介子 LCDA、$\Delta G$、双部分子分布 \tagP]\label{prop:lamet-other-objects}
（以下各项都只在单圈或在结构层次上建立。）(i) 重子 LCDA：准 DA 的单圈因子化“在单圈层次被显式证明”，给出 $\MSbar$ 与 RI/MOM 核\cite{Deng:2023csv}。(ii) 重介子 LCDA：两步因子化——QCD LCDA $=$ 微扰的类喷注函数 $\otimes$ HQET LCDA\cite{Beneke:2023nmj}；HQET LCDA 可通过带快速运动重夸克的等时关联函数由 LaMET 获得\cite{Han:2024fkr}，要求层次 $\Lambda\ll m_Q\ll P^z$。(iii) $\Delta G$ 是规范不变、依赖参考系的胶子自旋 $\vec E\times\vec A_\perp$ 的大动量极限\cite{Ji:2013fga}，其普适类只有 $\order{1/P_z^n}$（$n\ge2$）的修正\cite{Hatta:2013gta}。(iv) 双部分子分布：对所有味、自旋与色结构的 DPD 与准 DPD 的匹配，给出单圈夸克—夸克系数\cite{Jaarsma:2023woo}；色单态非极化夸克 DPD 的因子化公式\cite{Zhang:2023wea}（预印本）。以上均据其摘要。
（本命题只涉及部分子态与匹配系数。）
\uses{def:lattice-quasi, thm:lattice-lamet}
\end{proposition}

\paragraph{地位与缺口}
\tagP：单圈或结构性陈述；色非单态 DPD 的全阶软/快度结构与多规范链重子算符的重整化都没有建立；重介子的两步链合起来是 \tagP（每一步单圈显式）。一般的断言——这些量都满足与\cref{thm:lattice-lamet}同类的全阶匹配——是 LaMET 原理（\cref{conj:lamet-principle}）的特例。
